\documentclass[11pt]{article}
\usepackage[margin=2.1cm]{geometry}
\usepackage{booktabs}
\usepackage{amsmath}
\usepackage{amssymb}
\usepackage{enumitem}
\usepackage{array}
\usepackage{listings}
\usepackage{xcolor}
\usepackage[T1]{fontenc}
\usepackage{lmodern}

\setlength{\parindent}{0pt}
\setlength{\parskip}{3.6pt}
\newcommand{\code}[1]{\texttt{#1}}
\newcommand{\hr}{\vspace{2pt}\hrule\vspace{6pt}}
\renewcommand{\arraystretch}{1.15}

\definecolor{good}{RGB}{0,110,60}
\definecolor{bad}{RGB}{150,30,30}
\newcommand{\yes}{\textcolor{good}{\textbf{Write this}}}
\newcommand{\no}{\textcolor{bad}{\textbf{Not this}}}

\begin{document}

\begin{center}
{\Large\bfseries EE\,362 --- Computer Organization Laboratory}\\[2pt]
{\large\bfseries How to Write the Lab Report}\\[4pt]
Fall 2026 \quad$\cdot$\quad Lab TA: Gandhi Torres
\end{center}
\hr

This guide was written \emph{after} reading your Lab 0 and Lab 1 reports. Every
example below is taken from real submissions this term. Nothing here is a matter of
taste --- each rule corresponds directly to points on the rubric.

\section*{1. Where the points actually are}

\begin{center}
\begin{tabular}{lrl}
\toprule
Component & Pts & What it is really asking \\
\midrule
Pre-lab & 5 & Did you arrive with a design in your head? \\
Compiles cleanly, warnings explained & 25 & \textbf{Show me the Quartus output.} \\
Testbench and vectors, coverage justified & 20 & \textbf{Why is this set of vectors enough?} \\
Simulation results match the specification & 20 & \textbf{Show me the waveform, and read it to me.} \\
Report chapter quality & 15 & Tables, derivations, structure. \\
In-lab demonstration & 15 & Can you explain your own design? \\
\bottomrule
\end{tabular}
\end{center}

\textbf{Notice that 65 of the 100 points are evidence, not prose.} The most common
way to lose marks this term was not being wrong --- it was being \emph{unsupported}.
Four reports in Lab 1 contained no figures at all; they were the four lowest scores,
and in at least one case the student clearly understood the material completely.

\begin{quote}
\textbf{The rule behind everything in this guide:} a claim without evidence is
worth nothing, and evidence without a sentence interpreting it is worth almost
nothing.
\end{quote}

\section*{2. The rule that costs the most marks: captions}

A figure caption must say \textbf{what the figure proves}. Not what it is.

\vspace{4pt}
\no
\begin{quote}
\emph{Figure 1: Plot from ModelSim.}\\
\emph{Figure 4: Majority of 4.}\\
\emph{Figure 1.4: tb\_voting4.}
\end{quote}

\yes\ --- all three of these are from your classmates' reports this term:
\begin{quote}
\emph{``With \code{En} low from 20\,ns to 40\,ns, \code{D} and \code{D8} change at
30\,ns while \code{Q} and \code{Q8} keep their values --- the latch is holding.''}

\emph{``\code{Y} is high only in the windows 30--40\,ns and 50--80\,ns, which are
rows 3, 5, 6 and 7 of Table 1.2.''}

\emph{``\code{AA} then \code{55} flip every bit, so a \code{generate} with a stuck
index would be visible here.''}
\end{quote}

The test: \textbf{cover the figure and read only the caption.} If the caption still
tells the reader what was proved, it is a good caption.

\textbf{Crop your screenshots.} A full-screen capture of Questa with the signals
two millimetres tall proves nothing. Crop to the interval that matters and, if your
tool allows it, put the cursors on it.

\section*{3. Tables: number the rows, then cite them}

\begin{center}
\begin{tabular}{c|ccc|c|l}
\toprule
\# & \code{Clr} & \code{En} & \code{D} & $Q^{+}$ & note \\
\midrule
0 & 1 & -- & -- & \textbf{0} & asynchronous clear, dominates \\
1 & 0 & 0 & -- & $\mathbf{Q}$ & hold --- the memory row \\
2 & 0 & 1 & 0 & \textbf{0} & transparent \\
3 & 0 & 1 & 1 & \textbf{1} & transparent \\
\bottomrule
\end{tabular}
\end{center}

Then write ``row 1'' in the prose, never ``the case where the enable is low and the
data does not matter''. Numbered rows let your Discussion, your coverage argument
and your figure captions all point at the same thing.

Every table and figure needs a number and a caption, and must be \textbf{referred to
by number} in the text. If nothing in your text says ``Table 2'', the table is
decoration.

\section*{4. Derivations: three lines, and name the law}

\begin{align*}
\text{(1) canonical} \quad & Y = \textstyle\sum m(3,5,6,7) = \bar{A}BC + A\bar{B}C + AB\bar{C} + ABC \\
\text{(2) simplified} \quad & Y = AB + AC + BC
   \qquad \textbf{[idempotent, distributive, complement laws]} \\
\text{(3) VHDL} \quad & \code{Y <= (A and B) or (A and C) or (B and C);}
\end{align*}

\no\ ``The equation simplifies to $AB + AC + BC$.''\\
\yes\ the three lines above, with the law named at step (2).

``It simplified'' is an assertion, not a derivation. About half of the Lab 0 reports
lost these marks, and recovering them takes one extra line of writing.

\section*{5. The coverage argument: the question nobody answered}

You cannot test everything. The rubric does not ask you to --- it asks you to
explain \textbf{why your vectors are enough}. The strongest form is:

\begin{quote}
\textbf{For each vector, name the wrong design it would catch.}
\end{quote}

\yes\ --- again, from real reports this term:
\begin{quote}
\emph{``All 8 rows are tested. The expected value counts the asserted inputs
instead of re-using $AB+AC+BC$, so a mistake in the algebra would not be hidden.''}

\emph{``Each select value is exercised with both values of the selected input
\textbf{and} both values of the unselected input, to prove the unselected input has
no influence.''}

\emph{``The disabled row is applied first, because a testbench that only drives
$E = 1$ cannot detect a decoder that ignores its enable.''}
\end{quote}

\textbf{A vector that cannot fail is not a test.} If two inputs carry the same
pattern, or every immediate is positive, or the latch is never closed while the data
changes, then the corresponding bug survives your entire testbench.

And make the testbench \textbf{self-checking}: \code{assert} statements, a counted
error total, and one \code{PASS}/\code{FAILED} line at the end. Looking at a
waveform and deciding it seems right only tests whether the output matches what you
already expected.

\section*{6. The Discussion: one real problem}

This section is required, it is graded, and the summary version loses marks.

\no\ \emph{``This lab went quite well. Each simulation matched the truth table
perfectly.''} \quad \emph{``The lab required only minor debugging.''}

\yes\ three sentences: \textbf{what broke}, \textbf{how you found it}, \textbf{what
it taught you}. Real examples from this term:

\begin{quote}
\emph{``The project could not find its own top level: Error (12007). The file
compiled, so the VHDL was not the problem --- the entity is \code{trans\_latch} with
an underscore and the top-level setting said \code{translatch}. Quartus matches the
name exactly.''}

\emph{``The first waveform looked correct, but all activity ended at 30\,ns.
Comparing the run against the requirements showed the testbench never proved clear
while holding or clear while transparent. I expanded it to eight checked vectors.''}

\emph{``The hold vector still passed, which made the bug harder to notice, because
the latch looked like it was holding when it was really just never updating.''}

\emph{``A missing \code{end if;} --- I used the reported line number and looked
slightly \emph{above} it, because the compiler reports where it gave up, not where
the mistake is.''}
\end{quote}

Every one of those is worth full marks. Note that none of them is about being
clever; they are about being specific.

\section*{7. The evidence set}

Each lab specification names the evidence to collect. Collect it \emph{during} the
session --- you cannot reconstruct a waveform at midnight.

\begin{center}
\begin{tabular}{ll>{\raggedright\arraybackslash}p{8.0cm}}
\toprule
E1 & Flow Summary + Messages & 0 errors; warnings sorted into fix / explain / ignore \\
E2 & RTL Viewer & did the tool understand what I wrote? \\
E7 & Waveform + transcript & does it behave, and can I prove it? \\
\bottomrule
\end{tabular}
\end{center}

\textbf{Read the numbers, do not just paste them.} ``Total registers: 0, as expected
for combinational logic'' is worth marks. A screenshot of the same number with no
sentence is not. In Lab 1 the right reading is the opposite: the latch report
\emph{should} list your latches, and 0 registers is still correct because a latch is
not a register.

\section*{8. Professional practice: the cheap marks}

\begin{itemize}[leftmargin=1.6em,itemsep=2pt]
\item \textbf{Your name on the title page.} Four reports this term had no name
anywhere. If I cannot tell whose work it is, I cannot grade it --- and I cannot give
you your pre-lab points either.
\item \textbf{A real filename:} \code{Lastname\_Lab1.pdf}. Not
\code{Notes\_260928\_220655.PDF}, not \code{Lab1 (1).pdf}.
\item \textbf{No placeholders.} \code{PASTE SCREENSHOT HERE}, \code{Insert value},
and \code{Placeholder: the Quartus RTL Viewer screenshot} all appeared in submitted
PDFs this term. That is an automatic deduction, and it is pure waste --- the work
was usually done.
\item \textbf{Delete the instructions.} If your report contains the words
``Recommended screenshot'' or a submission checklist, you submitted the worksheet
instead of the report.
\item \textbf{Code as text, not as a photograph.} Use
\code{\textbackslash lstinputlisting} so the code in the report is the code you
compiled. Pasted screenshots of code drift out of date the moment you fix a bug.
\item \textbf{Declare AI assistance} if you used it, as several of you did
correctly. Stating that a tool helped with document structure while the analysis is
your own is exactly the right professional habit.
\item \textbf{The template is a template, not a submission.} The model chapter that
ships in the Overleaf project is there to show the standard. Replace it with your
own work.
\end{itemize}

\section*{9. Before you submit --- tick every box}

\begin{center}
\begin{tabular}{>{\raggedright\arraybackslash}p{15.4cm}}
\toprule
$\square$\; My name is on the title page, and the filename has my name in it. \\
$\square$\; Every table has numbered rows and is cited by number in the text. \\
$\square$\; Every derivation has three lines and names its law. \\
$\square$\; \textbf{Every figure has one sentence saying what it proves.} \\
$\square$\; Screenshots are cropped to the part that matters and are legible. \\
$\square$\; The testbench is self-checking and the transcript is shown. \\
$\square$\; The coverage argument says which wrong design each vector catches. \\
$\square$\; E1, E2 and E7 are present, and each number is read in words. \\
$\square$\; The Discussion names one real problem, how I found it, what it taught. \\
$\square$\; No placeholders, no instruction text, no unfilled table cells. \\
$\square$\; The code in the report is included from the files I actually compiled. \\
\bottomrule
\end{tabular}
\end{center}

\vspace{6pt}
\hr
{\small \textbf{One sentence to remember.} A lab report is not a record of what you
did --- it is the \emph{argument} that your design is correct, and every figure,
table and number in it is there as evidence for that argument. Anything that is not
evidence is decoration, and decoration is worth zero.}

\end{document}
